\documentclass[12pt]{article}
\usepackage[a4paper,margin=1.8cm]{geometry}
\usepackage{amsmath,amssymb,booktabs,array}
\usepackage{xepersian}
\settextfont[Path=../1401-2/HW2 - official solution/,Extension=.ttf,BoldFont=*Bd,ItalicFont=*It,BoldItalicFont=*BdIt]{XB Niloofar}
\setdigitfont[Path=../1401-2/HW2 - official solution/,Extension=.ttf,BoldFont=*Bd,ItalicFont=*It,BoldItalicFont=*BdIt]{XB Niloofar}
\setlength{\parindent}{0pt}
\setlength{\parskip}{0.45em}
\begin{document}
\begin{center}{\Large\textbf{پاسخ رسمی تمرین ششم ـ نیم‌سال ۱۴۰۰--۱}}\end{center}

\section*{سؤال ۱}
با عملیات
\[
 R_2\leftarrow R_2-2R_1,\qquad
 R_3\leftarrow R_3-4R_1,\qquad
 R_4\leftarrow R_4-R_1
\]
داریم
\[
 R_2=[0,0,4,-3,0\mid-9],\qquad
 R_4=[0,0,4,-3,0\mid-5].
\]
پس \(R_4\leftarrow R_4-R_2\) سطر \([0,0,0,0,0\mid4]\) را می‌دهد؛ یعنی \(0=4\). بنابراین دستگاه \textbf{جواب ندارد}.

\section*{سؤال ۲}
روابط ژاکوبی
\[
 x_1^{(k+1)}=\frac{8-x_2^{(k)}+x_3^{(k)}}5,
 \quad x_2^{(k+1)}=\frac{x_1^{(k)}+x_3^{(k)}-7}{4},
 \quad x_3^{(k+1)}=\frac{1+x_1^{(k)}-x_2^{(k)}}4
\]
هستند. از \(X^{(0)}=0\):
\[
\begin{array}{c|rrr|c}
k&x_1&x_2&x_3&\|X^{(k)}-X^{(k-1)}\|_\infty\\ \hline
1&1.600000&-1.750000&0.250000&1.750000\\
2&2.000000&-1.287500&1.087500&0.837500\\
3&2.075000&-0.978125&1.071875&0.309375\\
4&2.010000&-0.963281&1.013281&0.065000\\
5&1.995313&-0.994180&0.993320&0.030898\\
6&1.997500&-1.002842&0.997373&0.008662\\
7&2.000043&-1.001282&1.000085&0.002712\\
8&2.000273&-0.999968&1.000331&0.001314\\
9&2.000060&-0.999849&1.000060&0.000271
\end{array}
\]
پس با معیار \(10^{-3}\) در گام نهم متوقف می‌شویم. جواب دقیق \(X=(2,-1,1)^T\) است.

\clearpage
\section*{سؤال ۳}
در روش توانی \(Y^{(k)}=AX^{(k-1)}\) را بر مؤلفهٔ با بزرگ‌ترین قدرمطلق تقسیم می‌کنیم:
\[
\begin{array}{c|r|rrr}
k&\mu_k&x_1&x_2&x_3\\ \hline
1&5.000000&1&-0.200000&-0.200000\\
2&2.600000&1&0.076923&0.076923\\
3&3.153846&1&-0.024390&-0.024390\\
4&2.951220&1&0.008264&0.008264\\
5&3.016529&1&-0.002740&-0.002740\\
6&2.994521&1&0.000915&0.000915\\
7&3.001830&1&-0.000305&-0.000305
\end{array}
\]
برای \(|\mu_k-\mu_{k-1}|<0.01\)، نخستین توقف گام هفتم است. زوج دقیق غالب \((3,(1,0,0)^T)\) است. اگر معیار تغییر بردار منظور باشد، گام ششم کافی است.

دایره‌های گرشگورین بازه‌های \([1,5]\)، \(\{-1\}\) و \([-3,-1]\) را می‌دهند؛ مقادیر ویژهٔ دقیق \(3,-1,-2\) هستند.

برای ماتریس \(B\)، \(B(1,1)^T=(2,2)^T\)، پس حدس اولیه دقیقاً بردار ویژهٔ مقدار \(2\) است و روش مقدار غالب \(5\) را نمی‌بیند. برای \(C\):
\[
 (1,1,1)^T\longmapsto(-1/3,-1/3,1)^T\longmapsto(1,1,1)^T.
\]
چرخهٔ دوگامی به علت مقادیر ویژهٔ \(3,3,-3\) و نبودن قدرمطلق غالب یکتا رخ می‌دهد.

\section*{سؤال ۴}
روابط گاوس--سایدل داده‌شده از \(X^{(0)}=(1,0,1)^T\) می‌دهند
\[
 \boxed{X^{(1)}=(-4,5.8,-21.32)^T},
\]
\[
 \boxed{X^{(2)}=(79.1866667,2.5546667,176.3808)^T}.
\]

\clearpage
\section*{سؤال ۵}
ماتریس تکرار ژاکوبی
\[
 T_J=\begin{bmatrix}0&-1/2&1/2\\1&0&1\\-1/2&-1/2&0\end{bmatrix}
\]
چندجمله‌ای مشخصهٔ \(\lambda(4\lambda^2+5)/4\) دارد، بنابراین
\[
 \rho(T_J)=\frac{\sqrt5}{2}>1.
\]
برای گاوس--سایدل
\[
 T_{GS}=\begin{bmatrix}0&-1/2&1/2\\0&-1/2&3/2\\0&1/2&-1\end{bmatrix}
\]
و مقادیر ویژه \(0,(-3+\sqrt{13})/4,(-3-\sqrt{13})/4\) هستند؛ پس
\[
 \rho(T_{GS})=\frac{3+\sqrt{13}}4>1.
\]
هر دو روش واگرا هستند. نبودن غلبهٔ قطری به‌تنهایی اثبات واگرایی نیست؛ دلیل قطعی در این‌جا شعاع طیفی است.

\section*{سؤال ۶}
پیاده‌سازی کامل دولیتل، کروت و چولسکی بدون تابع آماده در فایل
\lr{\texttt{HW6 - official solution.py}}
قرار دارد. برای نمونهٔ سؤال در روش کروت:
\[
 L=\begin{bmatrix}8&0&0\\-6&5/2&0\\2&-5/2&0\end{bmatrix},
 \qquad
 U=\begin{bmatrix}1&-3/4&1/4\\0&1&-1\\0&0&1\end{bmatrix}.
\]
ضرب مستقیم \(LU=A\) را تأیید می‌کند. ماتریس نمونه تکین است و \(\det A=0\)؛ صفر نهایی قطر \(L\) اشکال محاسباتی ایجاد نمی‌کند. این نمونه برای چولسکی مناسب نیست، چون مثبت‌معین نیست.
\end{document}
